At its core, AD is the recursive application of the chain rule:
%
\begin{equation}
	\frac{d}{dx}\left[f\left(g\left(x\right)\right)\right] = f^{\prime}\left(g\left(x\right)\right)g^{\prime}\left(x\right)
	\label{eq:chain_rule}
\end{equation}
%
To apply the chain rule to computer programs, we must reinterpret the sequence of instructions that a program performs as function composition.
For a complicated program with many steps, this is equivalent to nested composition $f(g(h(w(...(x)))))$.
For deep nesting, it may be more useful to imagine the flow of values through the program as a computational graph.
For the case of $f(g(x))$, there are three ``primal'' values: $x$, $g(x)$, and $f(g(x))$.
Correspondingly, there are two ``tangent'' values, $g^{\prime}(x)$ and $f^{\prime}(g(x))$.
This graph is shown in \autoref{fig:simple_ast} with square nodes indicating start and stop values, circular nodes indicating computation, and edges showing the flow of values.

\begin{figure}[ht!]
	\centering
	\begin{tikzpicture}[scale=1]
		\node (x) [value] at (0,0) {$x$};
		\node (g) [computation] at (0,-1.5) {$g$};
		\node (f) [computation] at (0,-3) {$f$};
		\node (v) [value] at (0,-4.5) {$f(g(x))$};
		\draw [->] (x) to node[left] {$t_{1}=1$} node[right] {$p_{1}=x$} (g);
		\draw [->] (g) to node[left] {$t_{2}=g^{\prime}(p_{1})t_{1}$} node[right] {$p_{2}=g(p_{1})$} (f);
		\draw [->] (f) to node[left] {$t_{3}=f^{\prime}(p_{2})t_{2}=f^{\prime}(g(x))g^{\prime}(x)$} node[right] {$p_{3}=f(p_{2})=f(g(x))$} (v);
	\end{tikzpicture}
	\caption{Computational graph of forward-mode evaluation of $f(g(x))$ }
	\label{fig:simple_ast}
\end{figure}

The values flowing down the right side of the graph in \autoref{fig:simple_ast} represent the normal evaluation of the program.
The values on the left represent the accumulated tangents, with the initial value ``seeded'' as unity.
Each primal computation relies on the values that immediately precede it.
Similarly, the tangent components rely on the previous tangent and primal component.
At the end of this graph, we have computed both the primal result and the tangent.
This accumulation of tangent values alongside the primal values to build the resulting derivative is known as forward-mode AD.

To extend forward-mode AD to multivariate functions, we set the tangent ``seed'' value of the independent variable of interest to unity, with the remaining variables' tangents set to zero.
Consider this more complicated computer program:
% 
\begin{algorithm}[H]
\caption{Example differentiable program}
\begin{algorithmic}
\STATE 
\STATE {$f$}$(a, b)$:
\STATE \hspace{0.5cm}$x \gets a^2 $
\STATE \hspace{0.5cm}$y \gets \sin{b}$
\STATE \hspace{0.5cm}\textbf{return}  $x*y$
\end{algorithmic}
\end{algorithm}
%
Suppose we want to find $\partial f/\partial b$.
The graph of all operations, primals, and tangents to do so is shown in \autoref{fig:ast}.

\begin{figure}[ht!]
	\centering
	\begin{tikzpicture}[scale=1]
		\node (a) [value] at (0,0) {$a$};
		\node (b) [value] at (2,0) {$b$};
		\node (pow) [computation] at (0,-2) {$x^{2}$};
		\node (sin) [computation] at (2,-2) {$\sin$};
		\draw [->] (a) to node[align=left,left] {$t_{1}=0$\\$p_{1}=a$} (pow);
		\draw [->] (b) to node[align=left,right] {$t_{2}=1$\\$p_{2}=b$}(sin);
		\node (times) [computation] at (1,-4) {$*$};
		\draw [->] (pow) to node[align=left,left,xshift=-5pt] {$t_{3}=2p_{1}t_{1}$ \\ $p_3=p_1^2$} (times);
		\draw [->] (sin) to node[align=left,right,xshift=5pt] {$t_{4}=\cos(p_{2})t_{2}$ \\ $p_{4}=\sin(p_{2})$}(times);
		\node (f) [value] at (1,-6) {$f(a,b)$};
		\draw [->] (times) to node[align=right, left] { $t_{5}=p_{3}t_{4} + t_{3}p_{4}$  $=a^2\cos(b)$}
		node[align=right, right] {$p_{5}=p_{3}*p_{4}$  $=a^2\sin(b)$}(f);
	\end{tikzpicture}
	\caption{Computational graph of forward-mode evaluation of $f(a,b)$, with the tangent seeded on $b$}
	\label{fig:ast}
\end{figure}

To compute $\partial f/\partial a$, we would set $t_{1}$ to unity and $t_{2}$ to zero and perform an additional sweep.
This generalizes to an approach for evaluating gradients for functions of the form $f: \R^{n} \rightarrow \R^{m}$.
Requiring one pass per independent variable implies that gradient evaluation in forward-mode has a complexity of $\mathcal{O}(n)$ for an $n$-dimensional input.
As such, forward-mode performs well for $f: \R \rightarrow \R^{n}$ but has no better complexity than finite-differencing for $f: \R^{n} \rightarrow \R$ (while maintaining machine-precision accuracy).

While forward-mode is simple and is straightforward to implement, the linear complexity poses a problem for high-dimensional functions such as large machine learning models or highly parameterized circuit optimizations.
A solution is to change the manner in which we decompose the chain rule.
Instead of computing the full derivative of each nested function recursively, we compute the derivate \textit{with respect to} each nested function recursively and work backwards.
In practice, this amounts to evaluating the primal trace of the function as before, but now collecting all primal values and their partial derivatives with respect to each parent value.
Then, we seed the ``adjoint'' of the dependent variable (output) with unity and work backwards through the graph, applying the chain rule utilizing the precomputed partials.
This is known as reverse-mode AD.

Looking again at our unary $f: \R \rightarrow \R$ function of $y = f(g(x))$, we draw the computational graph shown in \autoref{fig:reverse_simple_ast}.

\begin{figure}[ht!]
	\centering
	\begin{tikzpicture}[scale=1]
		\node (x) [value] at (0,0) {$x$};
		\node (g) [computation] at (0,-1.5) {$g$};
		\node (f) [computation] at (0,-3) {$f$};
		\node (y) [value] at (0,-4.5) {$y$};
		\draw [->, bend right] (x) to node[left] {$p_{1}=x$} (g);
		\draw [->, bend right] (g) to node[right] {$a_{1}=a_{2}\frac{\partial p_{2}}{\partial p_{1}}=a_{2}g^{\prime}(p_{1})=f^{\prime}(g(x))g^{\prime}(x)$} (x);
		\draw [->, bend right] (g) to node[left] {$p_{2}=g\left(p_{1}\right)$} (f);
		\draw [->, bend right] (f) to node[right] {$a_{2}=a_{3}\frac{\partial p_{3}}{\partial p_{2}}=a_{3}f^{\prime}(p_{2})$} (g);
		\draw [->, bend right] (f) to node[left] {$p_{3}=f\left(p_{2}\right)$} (y);
		\draw [->, bend right] (y) to node[right] {$a_{3}=1$} (f);
	\end{tikzpicture}
	\caption{Graph of reverse-mode evaluation of $f(g(x))$ }
	\label{fig:reverse_simple_ast}
\end{figure}

Here we again see the primal, normal program execution of the function on the left, and then the reverse pass accumulation of the adjoints on the right.
The resulting adjoint, $a_{1}$, is identical to the resulting forward-mode tangent $t_{3}$ from \autoref{fig:simple_ast}.
The difference is that in this case, all the intermediate values and the partial derivatives need to be stored, whereas in forward-mode, each step relies solely on the previous tangent and primal values in the graph.

The procedure becomes more complicated with less trivial, multi-arity functions.
Consider the function example from \autoref{fig:ast}, shown in \autoref{fig:reverse_ast} now in reverse-mode.

\begin{figure}[t]
	\centering
	\subfloat[Forward Sweep] {
		\begin{tikzpicture}[scale=1]
			\tikzstyle{every node}=[font=\small]
			\node (a) [value] at (0.5,0) {$a$};
			\node (b) [value] at (1.5,0) {$b$};
			\node (pow) [computation] at (0.5,-2) {$x^{2}$};
			\node (sin) [computation] at (1.5,-2) {$\sin$};
			\draw [->] (a) to node[align=left,left] {$p_{1}=a$} (pow);
			\draw [->] (b) to node[align=left,right] {$p_{2}=b$}(sin);
			\node (times) [computation] at (1,-4) {$*$};
			\draw [->] (pow) to node[align=left,left,xshift=0pt] {$p_3=p_1^2$ \\
				$\frac{\partial p_{3}}{\partial p_{1}} = 2p_{1}$
			} (times);
			\draw [->] (sin) to node[align=left,right,xshift=0pt,yshift=0pt] {$p_{4}=\sin(p_{2})$ \\
				$\frac{\partial p_{4}}{\partial p_{2}}=\cos(p_{2})$
			}(times);
			\node (f) [value] at (1,-6) {$f(a,b)$};
			\draw [->] (times) to node[left] {$p_{5}=p_{3}*p_{4}$}
			node[align=left,right] {$\frac{\partial p_{5}}{\partial p_{3}} = p_{4}$ \\ $\frac{\partial p_{5}}{\partial p_{4}} = p_{3}$}
			(f);
		\end{tikzpicture}
	}
    \quad
	\subfloat[Reverse Sweep] {
		\begin{tikzpicture}[scale=1]
			\tikzstyle{every node}=[font=\small]
			\node (a) [value] at (0.5,0) {$a$};
			\node (b) [value] at (1.5,0) {$b$};
			\node (pow) [computation] at (0.5,-2) {$x^{2}$};
			\node (sin) [computation] at (1.5,-2) {$\sin$};
			\draw [->] (pow) to node[align=left,left] {$a_{1}=a_{3}\frac{\partial p_{3}}{\partial p_{1}}$} (a);
			\draw [->] (sin) to node[align=left,right, xshift=0pt] {$a_{2}=a_{4}\frac{\partial p_{4}}{\partial p_{2}}$}(b);
			\node (times) [computation] at (1,-4) {$*$};
			\draw [->] (times) to node[align=left,left,xshift=0pt] {$a_{3} = a_{5}\frac{\partial p_{5}}{\partial p_{3}}$} (pow);
			\draw [->] (times) to node[align=left,right,xshift=0pt] {$a_{4}=a_{5}\frac{\partial p_{5}}{\partial p_{4}}$} (sin);
			\node (f) [value] at (1,-6) {$f(a,b)$};
			\draw [->] (f) to node[left] {$a_{5}=1$} (times);
		\end{tikzpicture}
	}
	\caption{Graph of reverse-mode evaluation of $f(a,b)$ }
	\label{fig:reverse_ast}
\end{figure}

Again, in the forward sweep, we collect primal values as well as partial derivates.
Then, in the reverse-mode, we collect adjoints down every branch by multiplying the previous adjoint by the appropriate precomputed  partial derivates.
This offers several advantages, the first of which is that in one pass, we have computed the full gradient.
Additionally, we have done so with fewer operations than in the forward-mode.
Often, this is a more efficient way of computing the gradient.
However, the reverse sweep requires us to build the computation graph in memory and store intermediate values and partials, while forward-mode only requires the previous tangents and primals.
This is a classic example of a memory-time complexity tradeoff.
In general, for functions of the form $f: \R^{n} \rightarrow \R^{m}$ where $m \gg n$, forward-mode is preferred.
However, for general optimization problems, reverse-mode is typically higher performance, as it solves $f: \R^{n} \rightarrow \R$ for a scalar cost function.
Benchmarking to choose the right approach is required when the input and output dimensions have similar orders of magnitude.

Efficient implementations of these algorithms have become commonplace due to their usage in training large machine learning models.
For a more in-depth discussion on the field of AD and implementation details, see~\cite{baydinAutomaticDifferentiationMachine2017}.